\documentclass[12pt,a4paper]{article}

% ─── Page Layout ───────────────────────────────────────────────────────────────
\usepackage[a4paper, margin=0.8in, headheight=14pt]{geometry}

% ─── Fonts (XeLaTeX) ───────────────────────────────────────────────────────────
\usepackage{fontspec}
\usepackage{unicode-math}
\setmainfont{TeX Gyre Pagella}
\setmonofont[Scale=0.88]{TeX Gyre Cursor}
\setmathfont{Latin Modern Math}

% ─── Packages ──────────────────────────────────────────────────────────────────
\usepackage{xcolor}
\usepackage{colortbl}
\usepackage{titlesec}
\usepackage{enumitem}
\usepackage{tabularx}
\usepackage{booktabs}
\usepackage{array}
\usepackage{amsmath}
\usepackage{tikz}
\usetikzlibrary{shapes.geometric, arrows.meta, positioning, calc, fit, decorations.pathreplacing, patterns}
\usepackage{tcolorbox}
\tcbuselibrary{skins, breakable}
\usepackage{hyperref}
\usepackage{fancyhdr}
\usepackage{graphicx}
\usepackage{microtype}
\hyphenpenalty=10000
\exhyphenpenalty=10000
\sloppy
\usepackage{caption}
\usepackage{setspace}
\usepackage{multirow}
\usepackage{tocloft}

% ─── Q&A helper ────────────────────────────────────────────────────────────────
\newcommand{\Q}[1]{\textbf{#1}\par\noindent\ignorespaces}

% ─── Colour Palette ────────────────────────────────────────────────────────────
\definecolor{myred}{RGB}{196,30,58}
\definecolor{mydark}{RGB}{30,30,50}
\definecolor{mygreen}{RGB}{14,120,80}
\definecolor{myteal}{RGB}{0,130,140}
\definecolor{myamber}{RGB}{180,100,0}
\definecolor{mypurple}{RGB}{100,30,160}
\definecolor{mygray}{RGB}{246,247,249}
\definecolor{mgframe}{RGB}{180,185,200}
\definecolor{watermark}{RGB}{145,150,170}
\definecolor{hdrred}{RGB}{250,232,235}
\definecolor{hdrgreen}{RGB}{220,244,233}
\definecolor{hdrteal}{RGB}{220,242,244}
\definecolor{hdramber}{RGB}{253,242,215}
\definecolor{hdrpurple}{RGB}{238,228,255}
\definecolor{hdrgray}{RGB}{238,240,245}
\definecolor{rowalt}{RGB}{250,251,253}

% ─── Section Formatting ────────────────────────────────────────────────────────
\titleformat{\section}
  {\large\bfseries\color{myred}}{\thesection.}{0.5em}{}
  [\vspace{1pt}{\color{myred}\rule{\linewidth}{0.8pt}}]
\titleformat{\subsection}
  {\normalsize\bfseries\color{myteal}}{\thesubsection}{0.5em}{}
\titlespacing*{\section}{0pt}{18pt}{7pt}
\titlespacing*{\subsection}{0pt}{11pt}{4pt}

% ─── Paragraph ─────────────────────────────────────────────────────────────────
\setlength{\parindent}{0pt}
\setlength{\parskip}{4pt}

% ─── TOC Styling ───────────────────────────────────────────────────────────────
\renewcommand{\cfttoctitlefont}{\large\bfseries\color{myred}}
\renewcommand{\cftaftertoctitle}{\par\noindent{\color{myred}\rule{\linewidth}{0.8pt}}}
\setlength{\cftbeforetoctitleskip}{0pt}
\setlength{\cftaftertoctitleskip}{6pt}
\renewcommand{\cftsecfont}{\normalsize\bfseries\color{mydark}}
\renewcommand{\cftsecpagefont}{\normalsize\bfseries\color{mydark}}
\renewcommand{\cftsecleader}{\cftdotfill{\cftdotsep}}
\setlength{\cftbeforesecskip}{7pt}
\renewcommand{\cftsubsecfont}{\small\color{mydark}}
\renewcommand{\cftsubsecpagefont}{\small\color{mydark}}
\setlength{\cftbeforesubsecskip}{3pt}
\setlength{\cftsubsecindent}{1.4em}

% ─── Table helpers ─────────────────────────────────────────────────────────────
\renewcommand{\arraystretch}{1.45}
\setlength{\tabcolsep}{6pt}
\arrayrulecolor{mydark}
\newcolumntype{B}[1]{>{\small\bfseries\raggedright\arraybackslash}m{#1}}
\newcolumntype{Y}{>{\small\raggedright\arraybackslash}X}
\renewcommand{\tabularxcolumn}[1]{m{#1}}

% ─── Header / Footer ───────────────────────────────────────────────────────────
\pagestyle{fancy}
\fancyhf{}
\fancyhead[L]{\small\color{myred}\textbf{ES-CS301}\;\color{mydark}--- Data Structure \& Algorithms}
\fancyhead[R]{\small\color{myteal}\textbf{Module 4 Notes}}
\fancyfoot[C]{\small\color{mydark}\thepage}
\fancyfoot[R]{\small\color{watermark}\textit{\copyright\ Biraj}}
\renewcommand{\headrulewidth}{0.5pt}
\renewcommand{\footrulewidth}{0.3pt}

% ─── Hyperref ──────────────────────────────────────────────────────────────────
\hypersetup{
  colorlinks=true, linkcolor=myteal, urlcolor=myteal,
  pdfauthor={Biraj Sarkar (CGEC)},
  pdftitle={ES-CS301 --- Module 4 Notes: Sorting, Hashing and Graph Theory}
}

% ─── tcolorbox styles ──────────────────────────────────────────────────────────
\tcbset{
  tikzbox/.style={
    colback=mygray, colframe=mgframe, boxrule=0.6pt, arc=4pt,
    left=8pt, right=8pt, top=6pt, bottom=6pt,
    colbacktitle=mgframe, coltitle=mydark,
    fonttitle=\small\bfseries
  },
  defbox/.style={
    colback=hdrgreen, colframe=mygreen, boxrule=0.8pt, arc=4pt,
    left=9pt, right=9pt, top=6pt, bottom=6pt,
    colbacktitle=mygreen, coltitle=white,
    fonttitle=\small\bfseries
  },
  formulabox/.style={
    colback=hdrteal, colframe=myteal, boxrule=0.8pt, arc=4pt,
    left=9pt, right=9pt, top=7pt, bottom=7pt,
    colbacktitle=myteal, coltitle=white,
    fonttitle=\small\bfseries
  },
  examplebox/.style={
    colback=hdramber, colframe=myamber, boxrule=0.8pt, arc=4pt,
    left=9pt, right=9pt, top=6pt, bottom=6pt,
    colbacktitle=myamber, coltitle=white,
    fonttitle=\small\bfseries
  },
  masonbox/.style={
    colback=hdrpurple, colframe=mypurple, boxrule=0.8pt, arc=4pt,
    left=9pt, right=9pt, top=6pt, bottom=6pt,
    colbacktitle=mypurple, coltitle=white,
    fonttitle=\small\bfseries
  },
  exambox/.style={
    colback=hdrpurple, colframe=mypurple, boxrule=1pt, arc=5pt,
    left=10pt, right=10pt, top=8pt, bottom=8pt,
    colbacktitle=mypurple, coltitle=white,
    fonttitle=\bfseries,
    breakable, title after break={\textbf{Frequently Asked / Viva Questions (contd.)}}
  }
}

\tcbset{breakable}
\tcbset{after={\color{black}}}

% ─── TikZ styles ───────────────────────────────────────────────────────────────
\tikzset{
  vnode/.style={circle, draw=myteal, fill=hdrteal, font=\small\bfseries, inner sep=2pt, minimum size=0.75cm},
  arrow/.style={-Stealth, thick, color=mydark},
  redarrow/.style={-Stealth, thick, color=myred},
  line/.style={thick, color=mydark}
}

\begin{document}

% ─── Title Page ───────────────────────────────────────────────────────────────
\begin{center}
  {\LARGE\bfseries\color{myred} ES-CS301 --- Data Structure \& Algorithms}\\[5pt]
  {\normalsize\color{mydark} Semester III \quad$\bullet$\quad 3L:0T:4P \quad$\bullet$\quad 3 Credits}\\[3pt]
  {\normalsize\color{myteal}\textbf{Module 4 Exam Notes: Sorting Algorithms, Hashing Techniques and Graph Theory}}\\[3pt]
  {\small\color{mydark} CGEC \;$|$\; B.Tech ECE (2023--27) \;$|$\; \textit{Biraj Sarkar}}
\end{center}
\vspace{2pt}
\noindent{\color{myred}\rule{\linewidth}{1.5pt}}
\vspace{2pt}

\hypersetup{linkcolor=mydark}
\tableofcontents
\hypersetup{linkcolor=myteal}
\newpage

% ───────────────────────────────────────────────────────────────────────────────
\section{Classification and Criteria of Sorting Algorithms}
% ───────────────────────────────────────────────────────────────────────────────

\begin{tcolorbox}[defbox, title={\textbf{Definition: Sorting Problem and Taxonomic Criteria}}]
Given a sequence of $n$ records $R_1, R_2, \dots, R_n$ with keys $K_1, K_2, \dots, K_n$, \textbf{sorting} is the algorithmic permutation of these records into an ordered sequence $R_{\pi(1)}, R_{\pi(2)}, \dots, R_{\pi(n)}$ such that:
\begin{equation}
  K_{\pi(1)} \le K_{\pi(2)} \le \dots \le K_{\pi(n)}
\end{equation}
\textbf{Core Taxonomic Criteria:}
\begin{enumerate}[leftmargin=*, nosep]
  \item \textbf{Internal vs External Sorting:} Internal sorting holds the entire dataset concurrently in main RAM. External sorting processes massive datasets residing on secondary storage (hard drives, tapes) using phased multi-way merge passes.
  \item \textbf{Space Efficiency (In-Place vs Out-of-Place):} An in-place algorithm requires strictly $\mathcal{O}(1)$ auxiliary space beyond the input buffer (e.g., Insertion Sort, Heap Sort). Out-of-place algorithms require $\mathcal{O}(n)$ auxiliary storage (e.g., Merge Sort).
  \item \textbf{Stability:} A sorting algorithm is \textbf{stable} if records with identical keys preserve their relative initial input ordering:
  \begin{equation}
    K_i = K_j \text{ with } i < j \implies \pi(i) < \pi(j)
  \end{equation}
  \item \textbf{Adaptivity:} An algorithm is adaptive if its runtime decreases when the input array is already partially or completely sorted (e.g., Bubble Sort with swap flag, Insertion Sort).
\end{enumerate}
\end{tcolorbox}

\begin{tcolorbox}[formulabox, title={\textbf{Theoretical Lower Bound for Comparison-Based Sorting}}]
Any deterministic comparison-based sorting algorithm can be modeled as a binary decision tree where each internal node represents a comparison $A[i] \le A[j]$ and each leaf represents one of the $n!$ possible permutations of the input.
\par\vspace{3pt}
For a binary tree of height $h$ with at least $n!$ leaves:
\begin{equation}
  2^h \ge n! \implies h \ge \log_2(n!)
\end{equation}
Applying Stirling's approximation ($\ln(n!) \approx n \ln n - n$):
\begin{equation}
  h \ge \sum_{k=1}^n \log_2 k \ge \sum_{k=n/2}^n \log_2\left(\frac{n}{2}\right) = \frac{n}{2} \log_2\left(\frac{n}{2}\right) = \mathbf{\Omega(n \log_2 n)}
\end{equation}
Thus, no comparison-based sorting algorithm can achieve a worst-case time complexity faster than $\mathbf{\Omega(n \log n)}$.
\end{tcolorbox}

% ───────────────────────────────────────────────────────────────────────────────
\section{Elementary Sorting Algorithms}
% ───────────────────────────────────────────────────────────────────────────────

\subsection{Selection Sort}
\begin{itemize}[leftmargin=*, itemsep=3pt]
  \item \textbf{Working Principle:} In each pass $i \in [0, n-2]$, find the minimum element in unsorted subarray $A[i \dots n-1]$ and swap it with $A[i]$.
  \item \textbf{Comparisons:} Total comparisons are fixed for all input distributions:
  \begin{equation}
    C(n) = \sum_{i=0}^{n-2} (n - 1 - i) = (n - 1) + (n - 2) + \dots + 1 = \frac{n(n - 1)}{2} = \mathbf{\Theta(n^2)}
  \end{equation}
  \item \textbf{Swaps:} At most $n - 1$ swaps ($\mathcal{O}(n)$ writes) --- ideal when write operations are physically expensive (e.g., EEPROM/Flash memory).
  \item \textbf{Properties:} In-place ($\mathcal{O}(1)$ space), \textbf{Unstable} (e.g., sorting $[4_a, 4_b, 2]$ swaps $4_a$ with $2$, placing $4_a$ after $4_b$).
\end{itemize}

\subsection{Bubble Sort}
\begin{itemize}[leftmargin=*, itemsep=3pt]
  \item \textbf{Working Principle:} Iteratively scan the array, comparing adjacent pairs $(A[j], A[j+1])$ and swapping them if out of order, bubbling the largest item to the end in each pass.
  \item \textbf{Optimization:} Maintain a boolean flag \texttt{swapped}. If no swaps occur in a full pass, the array is already sorted and the algorithm terminates immediately.
  \item \textbf{Complexity:} Best-case: $\mathbf{\mathcal{O}(n)}$ (with flag on sorted input). Average and Worst-case: $\mathbf{\mathcal{O}(n^2)}$ comparisons and swaps.
  \item \textbf{Properties:} In-place ($\mathcal{O}(1)$ space), \textbf{Stable}, Adaptive.
\end{itemize}

\subsection{Insertion Sort}
\begin{itemize}[leftmargin=*, itemsep=3pt]
  \item \textbf{Working Principle:} Simulates how cards are sorted in hand: incrementally pick $A[i]$ and insert it into its correct relative position within the sorted subarray $A[0 \dots i-1]$ by shifting larger elements one position to the right.
  \item \textbf{Complexity Analysis:}
  \begin{itemize}[leftmargin=*, nosep]
    \item \textbf{Best-Case (Sorted Input):} Only 1 comparison per element, 0 shifts $\implies \mathbf{\mathcal{O}(n)}$ time.
    \item \textbf{Worst-Case (Reverse-Sorted Input):} Max shifts: $\sum_{i=1}^{n-1} i = \frac{n(n-1)}{2} \implies \mathbf{\mathcal{O}(n^2)}$ time.
    \item \textbf{Average-Case:} $\frac{n(n-1)}{4} \text{ shifts} \implies \mathbf{\mathcal{O}(n^2)}$ time.
  \end{itemize}
  \item \textbf{Properties:} In-place ($\mathcal{O}(1)$ space), \textbf{Stable}, Online (can sort incoming streaming data), highly optimal for small arrays ($n \le 16$).
\end{itemize}

% ───────────────────────────────────────────────────────────────────────────────
\section{Advanced Divide-and-Conquer \& Tree-Based Sorting}
% ───────────────────────────────────────────────────────────────────────────────

\subsection{Quick Sort (Partition-Exchange Sort)}
\begin{tcolorbox}[defbox, title={\textbf{Quick Sort Principle and Partitioning Strategy}}]
Quick Sort operates on the \textbf{Divide-and-Conquer} paradigm:
\begin{enumerate}[leftmargin=*, nosep]
  \item \textbf{Partition:} Select a pivot element $P$. Rearrange array such that all elements $\le P$ precede $P$, and all elements $\ge P$ succeed $P$. The pivot $P$ settles into its final sorted position $q$.
  \item \textbf{Recurse:} Recursively apply Quick Sort to subarrays $A[\text{low} \dots q-1]$ and $A[q+1 \dots \text{high}]$.
\end{enumerate}
\end{tcolorbox}

\begin{tcolorbox}[formulabox, title={\textbf{Quick Sort Recurrence Relations and Asymptotic Analysis}}]
\textbf{1. Best-Case Analysis (Balanced Partition):}
Pivot consistently bisects the array into two equal halves of size $n/2$:
\begin{equation}
  T(n) = 2 T\left(\frac{n}{2}\right) + \Theta(n) \implies \mathbf{T(n) = \Theta(n \log_2 n)}
\end{equation}

\textbf{2. Worst-Case Analysis (Extremely Unbalanced Partition):}
Occurs when the array is already sorted or reverse-sorted and the first or last element is chosen as pivot, yielding subproblems of size $0$ and $n-1$:
\begin{equation}
  T(n) = T(n - 1) + T(0) + \Theta(n) = T(n - 1) + c n = c \sum_{i=1}^n i \implies \mathbf{T(n) = \Theta(n^2)}
\end{equation}

\textbf{3. Average-Case Analysis:} Assuming all permutations equally probable: $\mathbf{\mathcal{O}(n \log n)}$.
\par\vspace{3pt}\noindent
\textbf{Auxiliary Space:} $\mathcal{O}(\log n)$ stack frames for balanced recursion; degrades to $\mathcal{O}(n)$ in worst case. Quick Sort is \textbf{Unstable}.
\end{tcolorbox}

\subsection{Merge Sort}
\begin{tcolorbox}[formulabox, title={\textbf{Merge Sort Algorithm and Guaranteed Bounds}}]
Merge Sort recursively divides the array into two equal halves until single-element subarrays are reached, then merges them in sorted order using an auxiliary buffer:
\begin{equation}
  T(n) = 2 T\left(\frac{n}{2}\right) + \Theta(n) \quad \forall n > 1, \quad T(1) = 0
\end{equation}
By the Master Theorem (Case 2: $a = 2, b = 2, f(n) = \Theta(n) = \Theta(n^{\log_2 2})$):
\begin{equation}
  \mathbf{T(n) = \Theta(n \log_2 n) \quad \text{in Best, Average, and Worst Cases!}}
\end{equation}
\par\vspace{2pt}\noindent
\textbf{Key Trade-offs:}
\begin{itemize}[leftmargin=*, nosep]
  \item \textbf{Guaranteed Performance:} Never degrades to $\mathcal{O}(n^2)$.
  \item \textbf{Space Overhead:} Requires $\mathbf{\mathcal{O}(n)}$ auxiliary memory array.
  \item \textbf{Stability:} \textbf{Stable} (ideal for complex multi-attribute sorting).
  \item \textbf{Linked Lists:} Does not require random access; sorts linked lists in $\mathcal{O}(n \log n)$ time with $\mathcal{O}(1)$ extra space.
\end{itemize}
\end{tcolorbox}

\subsection{Heap Sort}
A \textbf{Binary Heap} is a Complete Binary Tree mapped sequentially into an array $A[0 \dots n-1]$ satisfying the Max-Heap property: $A[\text{parent}] \ge A[\text{child}]$.
\begin{itemize}[leftmargin=*, nosep]
  \item Array Index Mapping: $\text{Parent}(i) = \lfloor (i - 1)/2 \rfloor, \; \text{Left}(i) = 2i + 1, \; \text{Right}(i) = 2i + 2$.
\end{itemize}

\begin{tcolorbox}[defbox, title={\textbf{Heap Sort Procedures and Linear Build-Heap Proof}}]
\begin{enumerate}[leftmargin=*, nosep]
  \item \textbf{Max-Heapify ($O(\log n)$):} Fixes violation at root of a subtree assuming left and right children are valid heaps.
  \item \textbf{Build-Max-Heap ($O(n)$):} Calls \texttt{Max-Heapify} on all internal nodes from $\lfloor n/2 \rfloor - 1$ down to $0$:
  \begin{equation}
    \sum_{h=0}^{\lfloor \log_2 n \rfloor} \left\lceil \frac{n}{2^{h+1}} \right\rceil \mathcal{O}(h) = \mathcal{O}\left( n \sum_{h=0}^\infty \frac{h}{2^h} \right) = \mathcal{O}(n \cdot 2) = \mathbf{\mathcal{O}(n)}
  \end{equation}
  \item \textbf{Heap Sort Algorithm ($\mathcal{O}(n \log n)$):}
  \begin{enumerate}[leftmargin=*, nosep]
    \item Build a Max-Heap from input array: $\mathcal{O}(n)$.
    \item For $i = n - 1$ down to $1$:
    \begin{itemize}[leftmargin=*, nosep]
      \item Swap root $A[0]$ (current maximum) with $A[i]$.
      \item Reduce heap size by 1.
      \item Call \texttt{Max-Heapify}($A, 0$) on reduced heap: $\mathcal{O}(\log n)$.
    \end{itemize}
  \end{enumerate}
\end{enumerate}
\par\vspace{2pt}\noindent
\textbf{Complexity:} Best, Average, Worst time strictly $\mathbf{\mathcal{O}(n \log n)}$. In-place ($\mathbf{\mathcal{O}(1)}$ space). \textbf{Unstable}.
\end{tcolorbox}

\subsection{Comprehensive Sorting Comparison Matrix}
\par\vspace{6pt}\noindent
\begingroup
\setlength{\tabcolsep}{3.5pt}
\setlength{\extrarowheight}{2.5pt}
\renewcommand{\arraystretch}{1.25}
\begin{tabularx}{\linewidth}{|B{2.2cm}|B{1.8cm}|B{1.8cm}|B{1.8cm}|B{1.3cm}|B{1.3cm}|Y|}
\hline
\textbf{Algorithm} & \textbf{Best} & \textbf{Average} & \textbf{Worst} & \textbf{Space} & \textbf{Stable?} & \textbf{Primary Use Case} \\
\hline
Selection Sort & $\mathcal{O}(n^2)$ & $\mathcal{O}(n^2)$ & $\mathcal{O}(n^2)$ & $\mathcal{O}(1)$ & No & Minimal write operations. \\
\hline
Bubble Sort & $\mathcal{O}(n)$ & $\mathcal{O}(n^2)$ & $\mathcal{O}(n^2)$ & $\mathcal{O}(1)$ & Yes & Educational; nearly sorted data. \\
\hline
Insertion Sort & $\mathcal{O}(n)$ & $\mathcal{O}(n^2)$ & $\mathcal{O}(n^2)$ & $\mathcal{O}(1)$ & Yes & Small lists ($n \le 16$); streaming. \\
\hline
Quick Sort & $\mathcal{O}(n \log n)$ & $\mathcal{O}(n \log n)$ & $\mathcal{O}(n^2)$ & $\mathcal{O}(\log n)$ & No & Fastest general in-memory sort. \\
\hline
Merge Sort & $\mathcal{O}(n \log n)$ & $\mathcal{O}(n \log n)$ & $\mathcal{O}(n \log n)$ & $\mathcal{O}(n)$ & Yes & Linked lists; external sorting. \\
\hline
Heap Sort & $\mathcal{O}(n \log n)$ & $\mathcal{O}(n \log n)$ & $\mathcal{O}(n \log n)$ & $\mathcal{O}(1)$ & No & Guaranteed $\mathcal{O}(n \log n)$ in-place. \\
\hline
\end{tabularx}
\endgroup

% ───────────────────────────────────────────────────────────────────────────────
\section{Hashing Techniques and Collision Resolution}
% ───────────────────────────────────────────────────────────────────────────────

\begin{tcolorbox}[defbox, title={\textbf{Definition: Hashing, Hash Table and Load Factor}}]
\textbf{Hashing} is a transformation technique that maps large variable keys $K$ from a universe $U$ into fixed integer slot indices in a \textbf{Hash Table} $T[0 \dots m-1]$ via a mathematical \textbf{Hash Function} $h(K)$:
\begin{equation}
  h : U \longrightarrow \{0, 1, \dots, m - 1\}
\end{equation}
The \textbf{Load Factor ($\alpha$)} is the ratio of stored keys $n$ to table capacity $m$:
\begin{equation}
  \alpha = \frac{n}{m}
\end{equation}
A \textbf{Collision} occurs when two distinct keys hash to the same slot: $h(K_1) = h(K_2)$ for $K_1 \ne K_2$.
\end{tcolorbox}

\subsection{Standard Hash Functions}
\begin{enumerate}[leftmargin=*, itemsep=3pt]
  \item \textbf{Division Method:} $h(K) = K \bmod m$. The table size $m$ should be a prime number not close to a power of 2 or 10 to minimize clustering.
  \item \textbf{Mid-Square Method:} Square the key $K^2$, and extract $r$ middle digits/bits as the index. Table size $m = 2^r$.
  \item \textbf{Folding Method:} Partition $K$ into parts matching index digit length, then sum parts:
  \begin{itemize}[leftmargin=*, nosep]
    \item \textbf{Fold-Shift:} Add components directly: $h(K) = (p_1 + p_2 + \dots) \bmod m$.
    \item \textbf{Fold-Boundary:} Reverse the digits of the outer parts before adding.
  \end{itemize}
  \item \textbf{Multiplication Method:} $h(K) = \lfloor m (K A \bmod 1) \rfloor$, where $0 < A < 1$ (Knuth recommends $A \approx \frac{\sqrt{5}-1}{2} \approx 0.618033$).
\end{enumerate}

\subsection{Collision Resolution Techniques}
\begin{enumerate}[leftmargin=*, itemsep=4pt]
  \item \textbf{Separate Chaining (Open Hashing):} Each table slot is a header pointer to a linked list of records hashing to that slot.
  \begin{itemize}[leftmargin=*, nosep]
    \item Can accommodate $\alpha > 1$ (table never gets "full").
    \item Average search time: $\mathcal{O}(1 + \alpha)$.
  \end{itemize}
  \item \textbf{Open Addressing (Closed Hashing):} All records reside directly in the table array ($m \ge n \implies \alpha \le 1$). If collision occurs, search alternate slots via a systematic probe sequence $h(K, i)$:
  \begin{itemize}[leftmargin=*, itemsep=2pt]
    \item \textbf{Linear Probing:} $h(K, i) = (h'(K) + i) \bmod m$. Suffers from \textbf{Primary Clustering} (adjacent occupied slots merge into long chains, degrading search to $\mathcal{O}(n)$).
    \item \textbf{Quadratic Probing:} $h(K, i) = (h'(K) + c_1 i + c_2 i^2) \bmod m$. Eliminates primary clustering but suffers from \textbf{Secondary Clustering} (keys hashing to same initial slot trace identical probe paths).
    \item \textbf{Double Hashing:} $h(K, i) = (h_1(K) + i \cdot h_2(K)) \bmod m$. Probe offset depends on the key itself ($h_2(K)$). Completely eliminates both primary and secondary clustering! Condition: $\gcd(h_2(K), m) = 1$.
  \end{itemize}
\end{enumerate}

% ───────────────────────────────────────────────────────────────────────────────
\section{Graph Theory, Representations and Traversals}
% ───────────────────────────────────────────────────────────────────────────────

\begin{tcolorbox}[defbox, title={\textbf{Definition: Graph Data Structure and Core Terminologies}}]
A \textbf{Graph} $G = (V, E)$ consists of a non-empty finite set of vertices $V$ and a set of edges $E$ connecting pairs of vertices.
\begin{itemize}[leftmargin=*, itemsep=2pt]
  \item \textbf{Handshaking Lemma:} In any undirected graph:
  \begin{equation}
    \sum_{v \in V} \deg(v) = 2 |E|
  \end{equation}
  \textit{Corollary:} In any graph, the number of vertices with odd degree is always even.
  \item \textbf{Complete Graph $K_n$:} Every pair of distinct vertices is connected by a unique edge: $|E| = \frac{n(n-1)}{2}$.
  \item \textbf{Connected Component:} A maximal connected subgraph of an undirected graph.
  \item \textbf{Strongly Connected Digraph:} A directed graph where there exists a directed path between every ordered pair of vertices $(u, v)$.
\end{itemize}
\end{tcolorbox}

\subsection{Graph Storage Representations}
\par\vspace{4pt}\noindent
\begingroup
\setlength{\tabcolsep}{4.5pt}
\setlength{\extrarowheight}{3pt}
\renewcommand{\arraystretch}{1.3}
\begin{tabularx}{\linewidth}{|B{3.2cm}|B{2.8cm}|B{3.0cm}|Y|}
\hline
\textbf{Representation} & \textbf{Space Complexity} & \textbf{Edge Lookup $(u, v)$} & \textbf{Optimal Suitability} \\
\hline
\textbf{Adjacency Matrix} & $\mathcal{O}(V^2)$ & $\mathcal{O}(1)$ & Dense graphs ($|E| \approx |V|^2$); fast edge tests. \\
\hline
\textbf{Adjacency List} & $\mathcal{O}(V + E)$ & $\mathcal{O}(\deg(u))$ & Sparse graphs ($|E| \ll |V|^2$); DFS/BFS algorithms. \\
\hline
\textbf{Incidence Matrix} & $\mathcal{O}(V \cdot E)$ & $\mathcal{O}(E)$ & Theoretical network flows; hypergraphs. \\
\hline
\end{tabularx}
\endgroup

\subsection{Graph Search and Traversal Algorithms}
\begin{tcolorbox}[formulabox, title={\textbf{Breadth-First Search (BFS) and Depth-First Search (DFS)}}]
\textbf{1. Breadth-First Search (BFS):}
\begin{itemize}[leftmargin=*, nosep]
  \item Explores vertices in concentric layers outward from source $s$ using a \textbf{FIFO Queue}.
  \item Time Complexity: $\mathbf{\mathcal{O}(V + E)}$ with Adjacency List; $\mathcal{O}(V^2)$ with Adjacency Matrix. Space: $\mathcal{O}(V)$.
  \item Applications: Single-source shortest path in unweighted graphs, bipartite testing.
\end{itemize}

\par\vspace{3pt}\noindent
\textbf{2. Depth-First Search (DFS):}
\begin{itemize}[leftmargin=*, nosep]
  \item Explores deeply along each branch until a dead-end before backtracking, using a \textbf{LIFO Stack} or system recursion.
  \item Classifies edges into: \textbf{Tree edges}, \textbf{Back edges} (detects directed cycles!), \textbf{Forward edges}, and \textbf{Cross edges}.
  \item Time Complexity: $\mathbf{\mathcal{O}(V + E)}$. Space: $\mathcal{O}(V)$ recursion call stack.
  \item Applications: Topological Sorting of DAGs, Cycle detection, Connected components, Biconnected components.
\end{itemize}
\end{tcolorbox}

\newpage
% ===============================================================================
% MANDATORY QUICK REVISION PAGE
% ===============================================================================
\begin{center}
  {\large\bfseries\color{myred} Quick Revision --- Module 4}\\[3pt]
  {\small\color{mydark} Sorting, Hashing and Graph Theory $|$ ES-CS301}
\end{center}
\noindent{\color{myred}\rule{\linewidth}{1.2pt}}
\vspace{4pt}

\begin{tcolorbox}[formulabox, title={\textbf{Key Formulas, Recurrences and Complexities at a Glance}}]
\renewcommand{\arraystretch}{1.35}
\begin{tabularx}{\linewidth}{B{5.0cm} X}
Comparison Sort Lower Bound & $\Omega(n \log_2 n) \quad \text{via Decision Tree Height } h \ge \log_2(n!)$ \\[3pt]
Merge Sort Recurrence & $T(n) = 2T(n/2) + \Theta(n) \implies \Theta(n \log_2 n) \quad \text{(Guaranteed)}$ \\[3pt]
Quick Sort Recurrence (Avg) & $T(n) = 2T(n/2) + \Theta(n) \implies \Theta(n \log_2 n)$ \\[3pt]
Quick Sort Recurrence (Worst) & $T(n) = T(n-1) + \Theta(n) \implies \Theta(n^2)$ \\[3pt]
Linear Build-Heap Bound & $\sum_{h=0}^{\log n} \frac{n}{2^{h+1}} \mathcal{O}(h) = \mathcal{O}(n) \quad \text{(Linear Time)}$ \\[3pt]
Hash Table Load Factor & $\alpha = n / m$ \\[3pt]
Linear Probing Probe Offset & $h(K, i) = (h'(K) + i) \bmod m$ \\[3pt]
Double Hashing Probe Offset & $h(K, i) = (h_1(K) + i \cdot h_2(K)) \bmod m$ \\[3pt]
Graph Handshaking Lemma & $\sum_{v \in V} \deg(v) = 2 |E|$ \\[3pt]
Complete Graph $K_n$ Edges & $|E| = \frac{n(n-1)}{2}$ \\[3pt]
BFS / DFS Time Complexity & $\mathcal{O}(V + E) \quad \text{(with Adjacency List)}$ \\[3pt]
\end{tabularx}
\end{tcolorbox}

\begin{tcolorbox}[defbox, title={\textbf{One-Line Definitions}}]
\begin{itemize}[leftmargin=*, itemsep=2pt, topsep=0pt]
  \item \textbf{Stable Sort:} A sorting algorithm that preserves the relative order of duplicate elements.
  \item \textbf{In-Place Sort:} A sorting algorithm requiring strictly $\mathcal{O}(1)$ auxiliary space beyond the input buffer.
  \item \textbf{Max-Heap:} A complete binary tree where the key of every parent node is greater than or equal to its children.
  \item \textbf{Hashing:} A technique that transforms a key into an integer array index via a hash function.
  \item \textbf{Collision:} An event where two distinct keys evaluate to the identical hash table slot index.
  \item \textbf{Primary Clustering:} The tendency in linear probing for occupied slots to form long continuous blocks.
  \item \textbf{Double Hashing:} An open addressing technique where the probe step size is calculated by a second hash function.
  \item \textbf{Adjacency Matrix:} A $V \times V$ matrix indicating edge connectivity between vertex pairs.
  \item \textbf{Adjacency List:} An array of linked lists where each entry lists the direct neighbors of a vertex.
  \item \textbf{Breadth-First Search (BFS):} Concentric layer-by-layer graph traversal implemented via a FIFO queue.
  \item \textbf{Depth-First Search (DFS):} Branch-deep graph traversal implemented via a stack or recursion.
\end{itemize}
\end{tcolorbox}

\begin{tcolorbox}[exambox, title={\textbf{Frequently Asked / Viva Questions}}]
\begin{enumerate}[topsep=0pt, label=\textbf{Q\arabic*.}, itemsep=5pt]
  \item \Q{Why can comparison-based sorting algorithms not execute faster than $\mathcal{O}(n \log n)$ in the worst case?}
  Every comparison sort can be modeled as a decision tree with $n!$ leaves corresponding to all permutations. A binary tree of height $h$ has at most $2^h$ leaves. Therefore, $2^h \ge n! \implies h \ge \log_2(n!) \approx n \log_2 n - n \log_2 e = \Omega(n \log n)$. The height represents the worst-case number of comparisons.

  \item \Q{Under what conditions does Quick Sort degrade to its worst-case $\mathcal{O}(n^2)$ time?}
  When the partition is maximally unbalanced (e.g., one subproblem of size $n-1$ and one of size $0$). This occurs when the array is already sorted or reverse-sorted and the first or last element is consistently selected as the pivot.

  \item \Q{Why is Heap construction (\texttt{Build-Max-Heap}) $\mathcal{O}(n)$ rather than $\mathcal{O}(n \log n)$?}
  Although \texttt{Max-Heapify} takes $\mathcal{O}(h)$ time, most nodes in a binary tree reside near the leaves with very small heights $h$. Summing the cost across all heights $\sum_{h=0}^{\lfloor \log n \rfloor} \frac{n}{2^{h+1}} h = \frac{n}{2} \sum \frac{h}{2^h} = \frac{n}{2} \cdot 2 = \mathcal{O}(n)$.

  \item \Q{Compare Separate Chaining and Open Addressing for collision resolution.}
  Separate Chaining uses external linked lists for colliding elements; it never suffers from table overflow and handles load factor $\alpha > 1$. Open Addressing stores all elements directly within the array table ($\alpha \le 1$); it eliminates pointer memory overhead but requires collision probing and can suffer from clustering.

  \item \Q{What is the difference between Primary Clustering and Secondary Clustering?}
  Primary clustering occurs in linear probing when different keys hashing to adjacent slots merge into a single long contiguous chain. Secondary clustering occurs in quadratic probing when two distinct keys having the same initial hash index trace the exact same probe sequence.

  \item \Q{State the Handshaking Lemma and explain its primary corollary.}
  The Handshaking Lemma states that in any undirected graph, $\sum_{v \in V} \deg(v) = 2|E|$, because each edge contributes 1 degree to two incident vertices. The corollary is that the number of vertices with odd degree in any graph must be strictly even.

  \item \Q{Compare Adjacency Matrix and Adjacency List representations of a graph.}
  Adjacency Matrix uses $\mathcal{O}(V^2)$ space and checks edge existence in $\mathcal{O}(1)$ time; it is optimal for dense graphs ($|E| \approx |V|^2$). Adjacency List uses $\mathcal{O}(V + E)$ space and finds neighbors in $\mathcal{O}(\deg(v))$ time; it is optimal for sparse graphs ($|E| \ll |V|^2$).

  \item \Q{How does DFS detect cycles in a directed graph?}
  During DFS traversal, vertices are colored White (unvisited), Gray (in active recursion stack), and Black (completely processed). If an edge $(u, v)$ encounters a Gray vertex $v$, that edge is a \textbf{Back Edge}, which definitively indicates the presence of a directed cycle.

  \item \Q{Why is Merge Sort preferred over Quick Sort for sorting linked lists?}
  In linked lists, random access is not available, which makes Quick Sort partitioning slow. Merge Sort divides a linked list in $\mathcal{O}(n)$ using the slow-fast pointer technique, and merges two sorted lists in $\mathcal{O}(n)$ time and strictly $\mathcal{O}(1)$ auxiliary space without requiring contiguous arrays.

  \item \Q{What is Topological Sorting, and what type of graph is required?}
  Topological sorting is a linear ordering of vertices such that for every directed edge $(u, v)$, vertex $u$ appears before $v$. It is strictly defined only on a \textbf{Directed Acyclic Graph (DAG)}.
\end{enumerate}
\end{tcolorbox}

\begin{tcolorbox}[tikzbox, title={\textbf{Summary Matrix: Comprehensive Comparison of Sorting, Hashing and Graphs}}]
\par\vspace{3pt}\noindent
\begingroup
\setlength{\tabcolsep}{4.0pt}
\setlength{\extrarowheight}{2.5pt}
\renewcommand{\arraystretch}{1.25}
\begin{tabularx}{\linewidth}{|B{2.8cm}|B{3.2cm}|B{3.0cm}|Y|}
\hline
\textbf{Technique / Topic} & \textbf{Core Mechanism} & \textbf{Time / Space Complexity} & \textbf{Key Practical Application} \\
\hline
Quick Sort & Divide-and-Conquer;\\[1pt]Pivot partitioning & Avg: $\mathcal{O}(n \log n)$, Worst: $\mathcal{O}(n^2)$\\[1pt]Space: $\mathcal{O}(\log n)$ & Fastest general in-memory primitive data sorting. \\
\hline
Merge Sort & Divide-and-Conquer;\\[1pt]Two-finger merge & Best/Avg/Worst: $\mathcal{O}(n \log n)$\\[1pt]Space: $\mathcal{O}(n)$ & External storage sorting, linked list sorting. \\
\hline
Heap Sort & Binary Max-Heap;\\[1pt]Root-leaf exchange & Guaranteed $\mathcal{O}(n \log n)$\\[1pt]Space: $\mathcal{O}(1)$ in-place & Systems where $\mathcal{O}(n \log n)$ guarantee and low space are critical. \\
\hline
Double Hashing & Two hash functions;\\[1pt]$h_1(K) + i \cdot h_2(K)$ & Search: $\mathcal{O}(1)$ avg\\[1pt]Space: $\mathcal{O}(m)$ & Hash tables requiring zero primary/secondary clustering. \\
\hline
BFS & Level-order expansion;\\[1pt]FIFO Queue & $\mathcal{O}(V + E)$\\[1pt]Space: $\mathcal{O}(V)$ & Shortest paths in unweighted graphs, web crawlers. \\
\hline
DFS & Deep path exploration;\\[1pt]LIFO Stack / Recursion & $\mathcal{O}(V + E)$\\[1pt]Space: $\mathcal{O}(V)$ & Cycle detection, topological sorting, connected components. \\
\hline
\end{tabularx}
\endgroup
\end{tcolorbox}

\end{document}
